% =====================================================================
%  第 2 节（开头）与 2.1 节  增强的 Fukaya–Morse A∞ 范畴：定义（原书 p.3–5）
% =====================================================================
\section{增强的 Fukaya--Morse $\Ainf$ 范畴：构造与例子}\label{sec:2}

\begin{takeaway}
\begin{itemize}
  \item \textbf{物理图像}：把“对象到自身”的态射也算进来。$F_a$ 到 $F_b$（$a\neq b$）的态射是 Morse 链，$F_a$ 到自身的态射取为 $X$ 上的奇异链；高阶复合则数“经过给定一族链的 Morse 树”。
  \item \textbf{数学工具}：Morse 复形与余指标分级、奇异链复形及其余维数分级、Morse 树的模空间、横截性假设、循环 $\Ainf$ 复合。
  \item \textbf{主要结论}：得到的是一个\emph{部分定义}的 $\Ainf$ 范畴 $\mathbb{F}$——高阶复合只在满足横截性假设（假设 2.2）的输入上定义；式 \eqref{eq:7} 的 $\Ainf$ 关系式成立（定理~\ref{thm:1}）。
\end{itemize}
\end{takeaway}

本节描述 \cite{fukaya,fukayaoh,kontsevichsoibelman} 中 Fukaya--Morse $\Ainf$ 预范畴\footnote{\emph{预范畴}（$\Ainf$ pre-category）在 \cite[第 4.3 节]{kontsevichsoibelman} 的术语里，指高阶复合只对连接“横截”对象序列的态射序列才有定义（横截对象序列的类本身是 $\Ainf$ 预范畴数据的一部分）。\cite{fukayaoh} 中“拓扑 $\Ainf$ 范畴”的概念含有类似的保留：复合法只对对象序列的一个 Baire 子集有定义。我们对自己的增强不用这两个词，而说“部分定义”，原因是这里的横截性条件加在\emph{态射序列}上，而不只是加在态射所连接的对象上。}的一个增强版本。更确切地说，我们取 \cite[第 6.2 节]{kontsevichsoibelman} 所称“光滑函数的 Morse $\Ainf$ 预范畴”的一个全子范畴。增强之处在于给出对象到自身态射的一个模型。所得结构是一个部分定义的 $\Ainf$ 范畴：高阶复合映射只在一族稠密的输入态射序列上有定义。

\subsection{“增强”的 Fukaya--Morse $\Ainf$ 范畴的定义}\label{sec:2.1}

% [Tutor Note: 读本节前先记住一个总原则：这个范畴里所有“数出来的整数”都带符号，符号由 Morse 理论的有向计数给出；而所有横截性问题都靠“假设 2.2”回避，而不是靠扰动具体实现。作者把这一取舍称为“部分定义”，这是理解后文两个证明为何都用“无穷小邻域/光滑化”操作的前提。]

\begin{takeaway}
\begin{itemize}
  \item \textbf{物理图像}：态射 = 交点（$a\neq b$）或链（$a=b$）；复合 = 数树。树的每条边带双色 $(a,b)$ 且是 $F_a-F_b$ 的梯度轨道，三价顶点的颜色规则是 $(a,b),(b,c)\to(a,c)$。
  \item \textbf{数学工具}：$\Mor(F_a,F_b)=\mathrm{MC}_\bullet(F_a-F_b)$（按 Morse 余指标分级的自由 $\mathbb Z$-模）、$\Mor(F_a,F_a)=C^{\mathrm{sing}}_\bullet(X;\mathbb Z)$（按余维数分级）、模空间 $\mathcal M(\cdots)$ 的零维计数、循环配对。
  \item \textbf{主要结论}：复合映射 (3) 有 $n$ 个输入时上同调次数为 $2-n$；具循环性 (6) 的 $\Ainf$ 关系式 (7) 成立（定理~\ref{thm:1}）。
\end{itemize}
\end{takeaway}

\textbf{对象.}\quad 考虑 $\Ainf$ 范畴 $\mathbb{F}$（复合法部分定义），其对象集是闭 Riemann 流形 $X$ 上一族处于一般位置的光滑函数 $\{F_1,\dots,F_N\}$。特别地，我们假定差 $F_a-F_b$ 对 $a\neq b$ 都是 Morse--Smale 函数。\footnote{回忆一下：$X$ 上的函数 $F$ 称为 \emph{Morse 函数}，若其临界点是孤立的，且每个临界点处的 Hessian 非退化。若进一步对任意一对临界点 $P,Q$，$P$ 的稳定（上升）流形与 $Q$ 的不稳定（下降）流形横截相交，则称 $F$ 为 \emph{Morse--Smale 函数}。}

\textbf{态射.}\quad 对 $a\neq b$，由对象 $F_a$ 到对象 $F_b$ 的态射是 $F_a-F_b$ 的 Morse 链：
\begin{equation}\tag{1}\label{eq:1}
  \Mor(F_a,F_b):=\mathrm{MC}_\bullet(F_a-F_b)=\operatorname{span}_{\mathbb Z}\bigl(\mathrm{Crit}(F_a-F_b)\bigr),
\end{equation}
即 $F_a-F_b$ 的临界点的形式线性组合，其上同调次数取为临界点的 Morse 余指标。由 $F_a$ 到自身的态射空间，我们取为 $X$ 上光滑奇异链构成的复形：
\begin{equation}\tag{2}\label{eq:2}
  \Mor(F_a,F_a)=C^{\mathrm{sing}}_\bullet(X,\mathbb Z),
\end{equation}
其中链按余维数给出上同调次数。我们指出，自同态还有其它可能的模型，见第~\ref{sec:2.7} 节。

% [Tutor Note: 式 (1) 的下标约定极易出错：$P_{ab}$ 中前一个下标 $a$ 是源（$F_a$），后一个 $b$ 是靶（$F_b$），与函数差写成 $F_a-F_b$ 的顺序一致。若误按 $F_b-F_a$ 理解，Morse 指标会整体反号，式 (4) 的树形规则也随之颠倒。核对原书 p.4 式 (4) 与其后的三价顶点规则。]

% [Tutor Note: 提示“按余维数分级”与“按余指标分级”是同一件事的两种说法，但只差一个整体次数平移：$\operatorname{codim}=\dim X-\operatorname{coind}$。作者在 (1)(2) 中并列使用两种分级，正是为了让式 (3) 的次数公式 $2-n$ 在两块上都成立；这一处若统一成“按维数分级”，所有公式都要整体加 $\dim X$ 的平移，切勿随手统一。]

\textbf{复合映射.}\quad 有高阶复合映射
\begin{equation}\tag{3}\label{eq:3}
\begin{aligned}
m\colon\quad
&\Mor(F_{a_1},F_{a_1})^{\otimes k_1}\otimes\Mor(F_{a_1},F_{a_2})\otimes\Mor(F_{a_2},F_{a_2})^{\otimes k_2}\otimes\Mor(F_{a_2},F_{a_3})\otimes\cdots\\
&\otimes\Mor(F_{a_{r-1}},F_{a_r})\otimes\Mor(F_{a_r},F_{a_r})^{\otimes k_r}
\ \longrightarrow\ \Mor(F_{a_1},F_{a_r}),
\end{aligned}
\end{equation}
其定义如下。

\begin{enumerate}
  \item[(i)] 对 $r\ge 2$：
  \begin{equation}\tag{4}\label{eq:4}
  \begin{aligned}
  &m\bigl(\{Z_{\alpha,a_1}\}_{\alpha=1}^{k_1};[P_{a_1a_2}];\{Z_{\alpha,a_2}\}_{\alpha=1}^{k_2};[P_{a_2a_3}];\cdots;[P_{a_{r-1}a_r}];\{Z_{\alpha,a_r}\}_{\alpha=1}^{k_r}\bigr)\\
  &\quad=\sum_{P_{a_1a_r}\in\mathrm{Crit}(F_{a_1}-F_{a_r})}
  \#\mathcal M\bigl(\{Z_{\alpha,a_1}\}_{\alpha=1}^{k_1};P_{12};\{Z_{\alpha,a_2}\}_{\alpha=1}^{k_2};P_{a_2a_3};\\
  &\qquad\cdots;P_{a_{r-1}a_r};\{Z_{\alpha,a_r}\}_{\alpha=1}^{k_r};P_{a_1a_r}\bigr)\cdot[P_{a_1a_r}].
  \end{aligned}
  \end{equation}
  \item[(ii)] 对 $r=1$，定义 $m\colon\Mor(F_a,F_a)^{\otimes k}\to\Mor(F_a,F_a)$ 为
  \begin{equation}\tag{5}\label{eq:5}
  \begin{aligned}
  m(Z_1)&=\partial Z_1,\\
  m(Z_1,Z_2)&=Z_1\cap Z_2,\\
  m(Z_1,\dots,Z_k)&=0\quad(k\ge 3).
  \end{aligned}
  \end{equation}
\end{enumerate}

其中：

\begin{itemize}
  \item $P_{ab}$ 是 $F_a-F_b$ 的临界点；$[P_{ab}]$ 表示把该临界点看作 $\Mor(F_a,F_b)$ 中的一个向量。
  \item $Z_{\alpha,a}$ 是 $X$ 上的链。
  \item $\mathcal M(\cdots)$ 是“Morse 树”的模空间——$X$ 中的二叉有根嵌入树，其 1 价顶点（“叶”）在 $P_{a_ia_{i+1}}$ 处，根在 $P_{a_1a_r}$ 处。每条边带一个双指标 $(a,b)$（$a\neq b$），并且是 $F_a-F_b$ 的一条梯度轨道。一个 1 价顶点 $P_{a_ia_{i+1}}$ 邻接一条 $(a_i,a_{i+1})$-边，根邻接一条 $(a_1,a_r)$-边。在 3 价顶点处，入边为 $(a,b)$、$(b,c)$，出边为 $(a,c)$，对某三个标号 $a,b,c$ 成立。我们只考虑满足如下附加条件的 Morse 树：每条链 $Z_{\alpha,a}$ 都与树的一条 $(a,b)$-边或 $(b,a)$-边相交（对某个 $b$）。此外，给定 $i$ 时，链 $Z_{\alpha,a_i}$（作为 $m$ 的输入）的排列顺序，应与沿树中连接 $P_{a_{i-1}a_i}$ 到 $P_{a_ia_{i+1}}$ 的路径行进时与各链 $Z_{\alpha,a_i}$ 相遇的先后顺序一致（特别地，$Z_{\alpha,a_1}$ 沿从根到 $P_{a_1a_2}$ 的路径排列，$Z_{\alpha,a_r}$ 沿从 $P_{a_{r-1}a_r}$ 到根的路径排列）。
  \item $\#\mathcal M(\cdots)$ 是 Morse 树的数目，当该模空间零维时；否则取为零。
  \item 复合 (4) 在某些横截性假设下才有定义——即下面的假设 2.2——加在输入态射上，也就是加在临界点 $P_{ab}$ 与链 $Z_{\alpha,a}$ 上。同样，式 (5) 中的二元运算是链的相交，这里假定链 $Z_1,Z_2$ 横截相交。这些假设正是 $\mathbb{F}$ 中复合只是部分定义的原因。
  \item 微妙之处：式 (4) 隐含假定 $a_1\neq a_r$。$a_1=a_r$ 时的复合可由循环性得到，见下文。
\end{itemize}

有 $n$ 个输入的复合映射 (3)，其上同调次数为 $2-n$。

% [Tutor Note: 式 (4) 中 $\#\mathcal M$ 的参数写作 $P_{12}$，而式 (4) 左侧与求和处写作 $P_{a_1a_2}$。这是原文的记号不一致（疑为缩写或笔误），译文照录不改；读者把它读作 $P_{a_1a_2}$ 即可。]

% [Tutor Note: 三价顶点的颜色规则 $(a,b),(b,c)\to(a,c)$ 是说：沿树从叶到根走，标号必须“接力”，且每经过一个三价顶点就换一次色。这条规则把树的组合型与复合的输入顺序绑在一起，是全节最实质的约定；它同时也是“不需要四价顶点”的原因——四价顶点的两种分解在同调层面相消（见第 4.2 节 (a) 类边界层片的讨论）。]

\textbf{$\Mor$ 上的微分.}\quad 只有一个输入态射的复合映射就是态射空间上的微分。特别地，$m\colon\Mor(F_a,F_b)\to\Mor(F_a,F_b)$ 在 $a\neq b$ 时是 Morse 微分，在 $a=b$ 时按上面的定义就是链上的边界算子。

\textbf{对偶与循环性.}\quad 态射空间上有自然的配对
\begin{equation*}
  \langle\,,\rangle\colon\Mor(F_a,F_b)\otimes\Mor(F_b,F_a)\to\mathbb Z.
\end{equation*}
当 $a=b$ 时，它是两条链相交点数的（带符号）计数。\footnote{我们假定相交是横截的，因而该配对也只是部分定义的。}当 $a\neq b$ 时，置 $\langle[P],[P']\rangle=\delta_{PP'}$——这里用到 $F_a-F_b$ 的临界点同时是 $F_b-F_a$ 的临界点。于是可以说，范畴 $\mathbb{F}$ 对态射具有对偶对象：$\Mor(F_a,F_b)^{\vee}=\Mor(F_b,F_a)$。利用这一点可以定义循环 $\Ainf$ 复合
\begin{equation}\tag{6}\label{eq:6}
  c(x_0,\dots,x_p):=\langle x_0,\,m(x_1,\dots,x_p)\rangle\in\mathbb Z,
\end{equation}
其中 $\{x_i\}$ 是 $\mathbb{F}$ 中的态射，满足 $\mathrm{target}(x_i)=\mathrm{source}(x_{i+1})$，指标 $i$ 取值于 $\mathbb Z_{p+1}$。循环复合与式 (4) 中出现的数 $\#\mathcal M$ 一致，并且在输入的循环置换下不变。

\textbf{$a_1=a_r$ 的情形.}\quad 式 (4) 隐含假定 $a_1\neq a_r$。当 $a_1=a_r$ 时，可由循环性——即式 (6) 在输入的循环置换下的不变性——反推出复合映射 $m$。例如，映射 $m\colon\Mor(F_1,F_2)\otimes\Mor(F_2,F_1)\to\Mor(F_1,F_1)$ 把 $[P]\otimes[Q]$（$P,Q$ 是 $F=F_1-F_2$ 的两个临界点）送到 $\mathrm{Unstab}_P\cap\mathrm{Stab}_Q$，即 $P$ 的不稳定流形与 $Q$ 的稳定流形之交（相对于 Morse 函数 $F$）。

% [Tutor Note: 这里“反推”用的是循环形式的配对 $\langle x_0,m(x_1,\dots,x_p)\rangle$，而不是 (4) 的非循环版本；两者相差一个取对偶的手续，符号来自 $\dim X$ 次的交换（见前一条批注）。若直接对 (4) 取 $a_1=a_r$，会多出一个不属于树的边界项，这是最容易算错的一步。]

把式 (4) 限制在没有 $\Mor(F_a,F_a)$ 中输入态射（即没有链 $Z_{\alpha,a}$）的情形，就得到 \cite{fukaya,fukayaoh} 中通常的 Fukaya--Oh 构造。

\begin{thmenv}{定理}{1}\label{thm:1}
复合映射 $m$ 满足 $\Ainf$ 关系式：若 $x_1,\dots,x_n$ 是 $\mathbb{F}$ 中的一族态射，且对 $i=1,\dots,n-1$ 满足 $\mathrm{target}(x_i)=\mathrm{source}(x_{i+1})$，则
\begin{equation}\tag{7}\label{eq:7}
  \sum_{r\ge 0,\ s\ge 1,\ r+s\le n}
  \pm\,m\bigl(x_1,\dots,x_r,\ m(x_{r+1},\dots,x_{r+s}),\ x_{r+s+1},\dots,x_n\bigr)=0.
\end{equation}
\end{thmenv}

% [Tutor Note: 式 (7) 中的 $\pm$ 号在原文里没有显式写出，它是 Koszul 符号：当内层 $m$ 的次数为奇时要对换序引入符号，具体为 $(-1)^{|x_1|+\cdots+|x_r|}$ 型的因子。这是全节唯一“不显式”的地方，定理 1 的两个证明（第 3、4 节）都给出同一符号约定；若读者自行补证，务必先固定 Koszul 约定，否则会差出整体符号。]

% [Tutor Note: 定理 1 的证明不是本节给出，而是第 3、4 节的两个独立证明。这是本文的组织方式：先把范畴定义“数”出来，再用两种场论把关系式“证”出来。阅读时建议先接受定理 1，直接进入第 3 节看有效作用量如何把它变成一个 BV 主方程。]
\endinput
